
\chapter{Conclusão}

O presente trabalho teve como objetivo desenvolver um modelo de regressão linear múltipla para prever o consumo anual de energia elétrica residencial a partir de características dos domicílios, considerando fatores relacionados aos moradores, ao clima, ao imóvel e à quantidade de aparelhos eletrodomésticos.

Os resultados demonstraram que o modelo conseguiu identificar relações entre as variáveis preditoras e o consumo de energia elétrica. Entretanto, seu desempenho apresentou limitações, com coeficiente de determinação de 0,3625 no conjunto de teste, indicando que aproximadamente 36,25\% da variabilidade do consumo foi explicada pelas variáveis utilizadas.

A proximidade entre as métricas dos conjuntos de treino e teste, juntamente com o resultado da validação cruzada ($R^2 = 0,3523 \pm 0,0053$), indica que o modelo apresentou comportamento relativamente estável, sem evidências significativas de sobreajuste. Apesar disso, os valores de MAE e RMSE demonstram que ainda existem diferenças consideráveis entre os consumos previstos e observados.

Conclui-se que a regressão linear múltipla foi capaz de estabelecer uma relação entre as características residenciais e o consumo de energia elétrica, mas as variáveis selecionadas não foram suficientes para explicar grande parte de sua variação. Isso sugere que outros fatores não contemplados no conjunto de dados também podem influenciar o consumo, além de possíveis relações não lineares entre as variáveis.

Como possibilidade de trabalhos futuros, podem ser investigadas outras variáveis preditoras, técnicas de seleção de atributos e modelos de aprendizado de máquina capazes de representar relações não lineares, buscando melhorar a precisão das previsões.
